\chapter{Fundamentação Teórica}
\label{sec:fundamentacao}

\section{Segurança da Informação}
    Inicialmente a Segurança da Informação foi regida por três princípios:
    \begin{itemize}
        \item Confidencialidade: este princípio estabelece que uma determinada informação só pode 
        ser acessada por pessoas, sistemas ou entidades devidamente autorizadas, garantindo que seu
        conteúdo não seja divulgado ou exposto a terceiros sem permissão\citep{fontes2017segurancca}. 
        \item Integridade: assegura que os dados a serem transmitidos não devem, em hipotese alguma, ser alterado,
        corrompidos ou excluídos por terceiros durante o caminho do remetente até o destino\citep{fernandes2018segurancca}.    
        \item Disponibilida: determina que os recursos e informações de um sistema devem estar acessíveis aos usuários
        autorizados sempre que solicitados, minimizando interrupções e garantindo a continuidade das operações\citep{moreno2019introduccao}.
    \end{itemize}

    Com o avanço da tecnologia e a dinamização da internet apenas a tríade - CID - deixou de ser 
    suficinte para garantir a segurança das informações. Portanto houve a necessidade de incluir outros 
    três pilares: autencidade, responsabilidade e a confomidade legal.A autenticidade refere-se à garantia
    de que indivíduos e sistemas são realmente quem afirmam ser; a responsabilidade assegura rastreabilidade
    das ações; e a conformidade legal assegura aderência a legislações como LGPD e normas internacionais
    de segurança. \citep{iso27001_2013}


\section{Testes de Penetração (Pentest)}
    Os testes de penetração consistem em técnicas sistemáticas de identificação, exploração
    e validação de vulnerabilidades em sistemas computacionais. O pentest simula o comportamento 
    de um atacante real, permitindo avaliar o estado de segurança de uma organização de forma prática\citep{moreno2019introduccao}.

    Os testes podem ser classificados em:
    \begin{itemize}
        \item Caixa branca (white-box): o testador possui amplo conhecimento da infraestrutura;
        \item Caixa cinza (gray-box): conhecimento parcial;
        \item Caixa preta (black-box): sem conhecimento prévio do ambiente.
    \end{itemize}

    \subsubsection{Fases do Processo de Testes de Penetração}

        Os testes de penetração são conduzidos de forma estruturada, seguindo fases bem definidas que permitem
        a identificação, exploração e análise de vulnerabilidades de maneira sistemática. Metodologias consolidadas,
        como o Penetration Testing Execution Standard (PTES), o NIST SP 800-115 e o OWASP Testing Guide, adotam abordagens
        semelhantes, garantindo padronização e reprodutibilidade dos testes \citep{moreno2019introduccao}.

        
        

        \subsubsection*{Reconhecimento e Coleta de Informações}

        A fase de reconhecimento tem como objetivo coletar informações iniciais sobre o alvo do teste, como endereços IP,
         serviços expostos, tecnologias utilizadas e possíveis vetores de ataque. Essa etapa fornece a base para as fases
          subsequentes, direcionando as ações do testador e reduzindo incertezas no processo \citep{stallings2017network}.

        \subsubsection*{Enumeração de Serviços e Vulnerabilidades}

        Na fase de enumeração, aprofunda-se a análise dos serviços identificados, buscando versões de software, configurações
         e vulnerabilidades conhecidas. Ferramentas automatizadas auxiliam na correlação entre serviços detectados e falhas já
          documentadas, permitindo uma visão mais precisa da superfície de ataque \citep{moreno2019introduccao}.

        \subsubsection*{Exploração}

        A exploração consiste na tentativa controlada de explorar as vulnerabilidades identificadas, com o objetivo de 
        comprovar sua existência e avaliar o impacto real sobre o sistema. Essa fase é conduzida de forma ética e autorizada,
         respeitando o escopo previamente definido, sendo fundamental para diferenciar vulnerabilidades teóricas de falhas 
         efetivamente exploráveis \citep{stallings2017network}.

        \subsubsection*{Pós-exploração}

        A fase de pós-exploração busca avaliar as consequências do acesso obtido, analisando possibilidades de escalonamento
        de privilégios, movimentação lateral e acesso a informações sensíveis. Essa etapa permite mensurar os riscos reais
        associados às vulnerabilidades exploradas, relacionando-os aos princípios de confidencialidade, integridade e disponibilidade \citep{moreno2019introduccao}.

        \subsubsection*{Geração de Relatórios}

        A etapa final consiste na documentação dos resultados obtidos durante o teste de penetração. O relatório deve apresentar
        as vulnerabilidades identificadas, evidências técnicas, impactos potenciais e recomendações de mitigação. Essa fase é
        essencial para comunicar os achados de forma clara e subsidiar a tomada de decisões por gestores e equipes técnicas \citep{moreno2019introduccao}.

\section{Virtualização de Containers}
    A virtualização é amplamente utilizada para possibilitar a execução isolada de aplicações em um mesmo
    ambiente computacional. Tradicionalmente, esse isolamento é obtido por meio de máquinas virtuais, 
    que incluem sistemas operacionais completos, resultando em maior consumo de recursos. Como alternativa,
    a virtualização baseada em containers permite o compartilhamento do kernel do sistema operacional, 
    oferecendo isolamento em nível de processos de forma mais leve e eficiente \citep{pahl2015containerization}.

    O Docker consolidou-se como a principal plataforma de containers, possibilitando o empacotamento 
    de aplicações juntamente com suas dependências, garantindo portabilidade e consistência entre diferentes
    ambientes de execução \citep{merkel2014docker}.

    No contexto da segurança da informação, o uso de containers favorece a execução isolada de ferramentas,
    reduzindo conflitos de dependências e impactos no sistema hospedeiro. Quando corretamente configurados,
    containers oferecem isolamento adequado para aplicações sensíveis, sendo apropriados para ambientes de testes
    de segurança e laboratórios educacionais\citep{combe2016docker}.

    No contexto deste trabalho, o Docker é utilizado para virtualizar Nmap, Nuclei, Nikto, Gobuster, Dirsearch,
    Sqlmap, Commix, Caido, Netcat e Metasploit, garantindo execução uniforme e independente do sistema operacional
    do usuário.

\section{Ferramentas de Testes de Penetração}
    Ferramentas de testes de penetração são empregadas para apoiar a identificação, exploração e análise de vulnerabilidades
    em sistemas computacionais. Elas automatizam tarefas repetitivas, auxiliam na coleta de informações e permitem avaliar,
    de forma prática, a postura de segurança de redes, sistemas e aplicações. O uso dessas ferramentas
    é essencial para tornar o processo de pentest mais eficiente e sistemático, especialmente em ambientes complexos\citep{moreno2019introduccao}.

    A seguir, apresentam-se as principais ferramentas consideradas neste trabalho, classificadas de acordo com seu papel no processo de testes de penetração.

        \subsection{Nmap}

            O Nmap (Network Mapper) é uma ferramenta amplamente utilizada para reconhecimento e enumeração de redes. Seu principal
            objetivo é identificar hosts ativos, portas abertas, serviços em execução e possíveis sistemas operacionais. A fase de reconhecimento 
            é fundamental para o sucesso do pentest, pois fornece informações iniciais que orientam as etapas
            subsequentes de análise e exploração\citep{stallings2017network}.

        \subsection{Nuclei}

            O Nuclei é um framework de varredura baseado em templates, empregado para a detecção automatizada de vulnerabilidades conhecidas e más configurações.
            Seu modelo baseado em assinaturas facilita a identificação rápida de falhas recorrentes em aplicações e serviços, contribuindo para a eficiência da
            fase de enumeração de vulnerabilidades \citep{moreno2019introduccao}.
            
        \subsection{Nikto}

            O Nikto é um scanner de servidores web utilizado para identificar configurações inseguras, arquivos sensíveis expostos e vulnerabilidades conhecidas.
            Essa ferramenta auxilia na avaliação preliminar da segurança de aplicações web, permitindo detectar falhas 
            comuns que podem comprometer o ambiente\citep{moreno2019introduccao}.
        
        \subsection{Gobuster e Dirsearch}

            Ferramentas como Gobuster e Dirsearch são utilizadas para a enumeração de diretórios, arquivos e recursos ocultos em aplicações web.
            Essas ferramentas exploram listas de palavras para identificar caminhos não diretamente acessíveis, auxiliando na descoberta de superfícies
            de ataque adicionais. A enumeração detalhada de recursos expostos é uma etapa relevante na avaliação da 
            segurança de aplicações\citep{stallings2017network}.

        \subsection{SQLMap e Commix}
        
            O SQLMap é uma ferramenta destinada à identificação e exploração automatizada de vulnerabilidades de injeção SQL, enquanto o Commix
            foca na detecção de falhas de injeção de comandos no sistema operacional. Ambas reduzem o esforço manual do testador ao automatizar
            testes complexos, permitindo validar o impacto real das vulnerabilidades encontradas \citep{moreno2019introduccao}.

        \subsection{Metasploit}

            O Metasploit é uma plataforma abrangente para exploração de vulnerabilidades e realização de testes de pós-exploração.
            Ele permite avaliar as consequências práticas das falhas identificadas, auxiliando na compreensão do impacto sobre a
            confidencialidade, integridade e disponibilidade dos sistemas. Frameworks de exploração são 
            fundamentais para validar a efetividade dos mecanismos de segurança existentes\citep{stallings2017network}.

        \subsection{Netcat}

            O Netcat é uma ferramenta de rede de propósito geral utilizada para estabelecer conexões TCP/UDP, realizar testes de 
            conectividade e interagir com serviços de forma direta. Em atividades de pentest, é frequentemente empregado para validação
            de portas e serviços, simulação de comunicação cliente-servidor e apoio a testes que exigem troca de dados em baixo nível.
            Por permitir inspeção e interação com serviços de forma simples, o Netcat é útil em etapas de reconhecimento e validação de achados,
            complementando scanners automatizados e auxiliando na compreensão prática do comportamento de rede \citep{stallings2017network,moreno2019introduccao}.
        
        \subsection{Caido}

            O Caido é uma ferramenta utilizada como proxy de interceptação para análise de tráfego HTTP/HTTPS, permitindo a inspeção,
            modificação e repetição de requisições entre cliente e servidor. Esse tipo de ferramenta é amplamente empregado em testes
            de segurança de aplicações web, pois possibilita observar o funcionamento interno das comunicações e identificar falhas
            relacionadas a autenticação, autorização e validação de entradas. A análise do tráfego de aplicações web
            é fundamental para compreender o comportamento das requisições e detectar vulnerabilidades que não são facilmente identificadas
            por scanners automatizados \citep{moreno2019introduccao}.
            
            Ao atuar como intermediário entre o navegador e a aplicação alvo, o Caido permite que o testador avalie de forma prática
            o impacto de alterações em parâmetros e cabeçalhos HTTP, auxiliando na identificação de falhas lógicas e de configuração.
            Dessa forma, a utilização de ferramentas de interceptação complementa o uso de scanners e frameworks de exploração, 
            contribuindo para uma análise mais abrangente da segurança de aplicações web \citep{stallings2017network}.

\section{Interface de Gráfica}

    Interfaces gráficas (Graphical User Interface – GUI) contribuem para a usabilidade de ferramentas de segurança
    da informação ao reduzir a dependência do uso direto de comandos em linha de comando, especialmente para usuários
    iniciantes. Em um cenário de crescente complexidade dos ataques cibernéticos e das ferramentas de defesa, 
    soluções que abstraem detalhes operacionais favorecem a compreensão dos processos de testes de segurança e 
    reduzem erros de execução \citep{ahmad2023zero}.

    Plataformas educacionais voltadas ao pentest tornam-se mais eficazes quando disponibilizam interfaces interativas,
    permitindo maior foco no aprendizado prático dos conceitos de segurança. Nesse contexto, o uso de interfaces gráficas
    auxilia na redução da barreira de entrada para estudantes e profissionais em formação\citep{alaryani2024penthack}.

    No desenvolvimento em Python, o Tkinter apresenta-se como uma solução adequada por ser uma biblioteca nativa, simples
    e multiplataforma. A adoção do ttkbootstrap complementa essa abordagem ao proporcionar uma aparência visual mais
    moderna, mantendo a simplicidade da implementação. Assim, a interface gráfica adotada neste trabalho alinha-se aos
    objetivos educacionais da ferramenta proposta.

\section{Modelos de Linguagem de Grande Escala (LLMs) no Apoio ao Pentest}
    Modelos de Linguagem de Grande Escala (Large Language Models – LLMs) têm sido progressivamente incorporados a sistemas de apoio
    à segurança da informação, devido à sua capacidade de processar grandes volumes de dados textuais, interpretar resultados técnicos
    e fornecer explicações em linguagem natural. Esses modelos mostram-se particularmente úteis em ambientes complexos, nos quais a
    análise manual de informações pode ser custosa e sujeita a erros \citep{ahmad2023zero}.

    No contexto dos testes de penetração, LLMs podem atuar como ferramentas de suporte à interpretação de resultados, auxílio na compreensão
    de vulnerabilidades e orientação sobre boas práticas de segurança. O crescimento e a sofisticação dos
    ataques cibernéticos exigem abordagens que combinem automação e inteligência para apoiar profissionais e estudantes na tomada
    de decisão durante atividades de segurança \citep{ahmad2023zero}.

    Plataformas educacionais que integram inteligência artificial ao processo de aprendizado prático tendem a apresentar ganhos significativos
    na assimilação do conhecimento. O uso de IA em ambientes de pentest pode auxiliar usuários iniciantes
    ao fornecer explicações contextualizadas, sugestões de ações e suporte cognitivo durante a execução das ferramentas, reduzindo a dependência
    exclusiva de conhecimento prévio aprofundado \citep{alaryani2024penthack}.

    É importante destacar que, no contexto deste trabalho, o uso de LLMs não se destina à automação completa de ataques, mas sim ao apoio educacional
    e interpretativo, auxiliando o usuário na compreensão dos resultados obtidos e no correto uso das ferramentas de testes de penetração. Dessa forma,
    a integração de LLMs contribui para tornar o processo de aprendizagem mais acessível, alinhando-se aos objetivos formativos da ferramenta proposta.





    

    

    


    
